\documentclass[12pt,a4paper]{article}

%% ── Packages ────────────────────────────────────────────────────────────
\usepackage{amsmath,amssymb,amsfonts}
\usepackage{graphicx}
\usepackage[margin=2.5cm]{geometry}
\usepackage{booktabs}
\usepackage{multirow}
\usepackage{makecell}
\usepackage{array}
\usepackage{algorithm}
\usepackage{algpseudocode}
\usepackage[colorlinks=true,linkcolor=blue,citecolor=blue,urlcolor=blue,breaklinks=true]{hyperref}
\usepackage{url}
\usepackage{enumitem}
\usepackage{caption}
\usepackage[T1]{fontenc}
\usepackage{lmodern}
\usepackage{microtype}

%% ── Custom operators ────────────────────────────────────────────────────
\DeclareMathOperator*{\argmax}{argmax}
\DeclareMathOperator*{\argmin}{argmin}

%% ── Graphics path ───────────────────────────────────────────────────────
\graphicspath{{Q1_corrected_paper_revised_assets/}}

%% ── Convenience macros ──────────────────────────────────────────────────
\newcommand{\yes}{$\checkmark$}
\newcommand{\no}{$\times$}

%% ── Title ───────────────────────────────────────────────────────────────
\title{\textbf{SecFedRL: A Byzantine-Resilient Framework for\\
Decentralized Federated Reinforcement Learning}}
\author{}
\date{}

\begin{document}
\maketitle
\thispagestyle{empty}

%% ════════════════════════════════════════════════════════════════════════
%% ABSTRACT
%% ════════════════════════════════════════════════════════════════════════
\begin{abstract}
Federated Reinforcement Learning (FRL) allows multiple agents to collaboratively learn
a policy without exchanging their data. Though earlier models of FRL had centralized
servers for collecting data, recent Peer-to-Peer (P2P) FRL systems avoid Single Point of
Failure (SPOF) issues while introducing security concerns. P2P FRL systems are easily
attacked by Byzantine attackers via model updates, especially when they use stochastic
and non-IID distributions. Therefore, robust centralization techniques cannot perform
effectively against them because they have high variance gradients and no supervisors.
To address these limitations, SecFedRL proposes a flexible algorithm that is resistant
to Byzantine attacks. Contrary to conventional solutions in parameter spaces, SecFedRL
employs ground-truth-free policy-space behavior analyses and introduces proactive
collapse-prevention strategies for stable learning with adversarial agents. It identifies
compromised models using virtual trajectory analysis and seven consistency measures,
which include KL-divergence, cosine similarity, policy agreement, value alignment,
entropy divergence, reverse KL, and magnitude consistency. Moreover, the framework uses
reputation-based weightings and decentralized consensus approaches. In addition, the
system adopts a lightweight hash-chained audit ledger employing SHA-256 hash chaining
with two-thirds consensus providing tamper-evident verification for auditability purposes.
Experimentally, SecFedRL achieves a 98\% average per-round detection rate and 100\%
cumulative detection (True Positive Rate) against Byzantine attacks on the
CartPole-v1, LunarLander-v2, and Pendulum-v1 tasks with 10--50 agents and up to 30\%
Byzantines (95\% at 40\% Byzantines), attaining 88--97\% task success rates across
environments (CartPole: 97\%, LunarLander: 96\%, Pendulum: 88\%) and generating
362--521\% higher episodic rewards than the undefended FedAvg baseline across tested
environments (CartPole-v1: +521\%, LunarLander-v2: +362\%, Pendulum-v1: +78\% absolute
reward improvement), with a computational overhead of ${<}8.3\%$.
\end{abstract}

\noindent\textbf{Keywords:} Federated RL, Byzantine Resilience, P2P FRL, Secure
Aggregation, Distributed Consensus, Actor-Critic.

%% ════════════════════════════════════════════════════════════════════════
%% §2  INTRODUCTION
%% ════════════════════════════════════════════════════════════════════════
\setcounter{section}{1}
\section{Introduction}
\label{sec:intro}

Recent advancements in Deep Reinforcement Learning (DRL) have enabled the development
of highly autonomous and intelligent systems across a wide range of applications,
including automated warehouses, autonomous vehicles, UAV swarms, and IoT sensor
networks. These systems require advanced sequential decision-making capabilities. In
this context, recent studies have explored the integration of reinforcement
learning-based trust mechanisms into Byzantine-resilient Federated Learning (FL)
frameworks~\cite{ref1}. However, as these systems scale, significant concerns arise
regarding the privacy of trajectory data and the increasing centralization of learning
processes, which may introduce vulnerabilities and create single points of failure.
Kairouz et al.~\cite{ref2} identified robustness, privacy, and communication efficiency
as key open challenges in FL.

One of the major difficulties in FL is handling non-IID data distributions.
Zhu et al.~\cite{ref3} demonstrated that data heterogeneity can severely impact both
convergence speed and overall model performance. To mitigate privacy risks, Bonawitz
et al.~\cite{ref4} proposed a practical secure aggregation protocol for
privacy-preserving machine learning. Meanwhile, Shejwalkar et al.~\cite{ref5}
critically evaluated poisoning attacks in real-world FL systems, exposing limitations
and failure cases in traditional defense mechanisms.

Decentralized approaches to FL have also gained attention. Sokhankhosh and
Rouhani~\cite{ref6} introduced a blockchain-based multi-winner consensus mechanism for
collaborative learning, while Li and Lyu~\cite{ref7} provided convergence guarantees
for sequential FL under heterogeneous data settings. On the adversarial front,
Shejwalkar and Houmansadr~\cite{ref8} showed that model poisoning attacks can be
strategically optimized to maximize damage while remaining undetected.

In addition, prior research has explored several techniques to improve efficiency and
robustness in FL. These include communication-efficient methods for edge
devices~\cite{ref9}, layer-adaptive sparsified aggregation for Byzantine
resilience~\cite{ref10}, ensemble distillation for combining heterogeneous
models~\cite{ref11}, gradient quantization for reducing communication
overhead~\cite{ref12}, and over-the-air aggregation for wireless
networks~\cite{ref13}. Deressa and Hasan~\cite{ref14} proposed TrustBandit, a
reinforcement learning-based approach for optimizing client selection in the presence
of poisoning attacks.

Building on these developments, FRL has emerged as a promising paradigm for
collaborative policy learning. Fan et al.~\cite{ref15} presented the first theoretical
framework for fault-tolerant FRL. Fang et al.~\cite{ref16} demonstrated that carefully
designed poisoning attacks can bypass existing Byzantine-robust aggregation techniques.
Xie and Song~\cite{ref17} addressed data heterogeneity in FRL using KL-divergence
regularization, while Ni et al.~\cite{ref18} proposed rFedFW, a reputation-based
aggregation framework designed for Byzantine-robust learning in IoT environments.

The vulnerability of federated systems becomes even more pronounced in reinforcement
learning due to several inherent characteristics:
\begin{enumerate}[label=(\roman*)]
  \item stochastic exploration leads to high variance in policy gradients,
  \item non-IID trajectory distributions make it difficult to cluster benign updates,
  \item both actor and critic networks are exposed to attacks, and
  \item temporal error propagation arises from bootstrapped value
        estimation~\cite{ref15,ref17}.
\end{enumerate}

Under such conditions, traditional aggregation methods---including Krum~\cite{ref19},
geometric median approaches~\cite{ref20}, generalized Byzantine-tolerant
SGD~\cite{ref21}, and coordinate-wise trimmed mean~\cite{ref22}---often fail, as they
rely on the presence of dense clusters of benign updates. Lyu et al.~\cite{ref23}
further provided a comprehensive taxonomy of threats, including data poisoning, model
poisoning, and inference attacks.

FL, formally introduced by Yang et al.~\cite{ref24} and operationalized through the
FedAvg algorithm by McMahan et al.~\cite{ref25}, enables multiple distributed
participants to collaboratively train models by sharing parameters instead of raw data.
Recent advancements, such as cryptographic verification techniques proposed by Cui
et al.~\cite{ref26} and robust aggregation methods like FLRAM~\cite{ref27}, have
strengthened the overall defense landscape. Hou et al.~\cite{ref28} introduced a
decentralized Byzantine-robust framework based on feature extraction layers, while Xia
et al.~\cite{ref29} demonstrated the effectiveness of distributed optimization
techniques for achieving Byzantine resilience. Additionally, Salgia and
Chi~\cite{ref30} explored the trade-offs between sample complexity and communication
efficiency in federated Q-learning.

Despite these significant advancements, designing Byzantine-resilient FRL systems that
effectively combine behavioral analysis, reputation-based aggregation, and stability
mechanisms under adversarial and heterogeneous conditions remains an open research
challenge.

To address this gap, we propose SecFedRL, a fully decentralized Byzantine-resilient
framework operating over a P2P mesh topology for FRL. Unlike traditional
parameter-space defenses that rely on gradient
filtering~\cite{ref19,ref20,ref21,ref22,ref27} or sparsification
techniques~\cite{ref10}, SecFedRL operates entirely in the policy space by generating
virtual trajectories to analyze agent behavior without any central coordinator. It
computes multi-dimensional consistency scores across peer updates, allowing for more
accurate identification and filtering of malicious participants prior to aggregation.
The framework also integrates reputation-based weighting inspired by trust-driven
approaches~\cite{ref1,ref14,ref18}, incorporates two-thirds consensus validation
mechanisms derived from blockchain-based systems~\cite{ref6,ref26}, and includes
collapse-prevention strategies to maintain stable learning under adversarial conditions.

Extensive experimental evaluations on CartPole-v1, LunarLander-v2, and Pendulum-v1
demonstrate that SecFedRL achieves a 98\% detection rate (True Positive Rate) against
Byzantine adversaries. It also delivers 362--521\% improvement in cumulative rewards
over the undefended baseline across environments (CartPole-v1: +521\%,
LunarLander-v2: +362\%) and maintains 88--97\% task success rates (CartPole: 97\%,
LunarLander: 96\%, Pendulum: 88\%) under 20--30\% adversarial participation, with only
a modest computational overhead of 4.9--8.3\%.

\subsection*{Main Contributions}

The \textbf{primary technical contribution} of SecFedRL is a \textit{ground-truth-free
policy-space Byzantine detection} mechanism: rather than comparing raw gradient or
parameter vectors (as all prior Byzantine defenses do), SecFedRL evaluates peers by
running their models on shared synthetic states from each agent's replay buffer and
measuring seven behavioral consistency metrics. This policy-space approach is
fundamentally different from parameter-space defenses (Krum~\cite{ref19}, Trimmed
Mean~\cite{ref22}) and does not require knowledge of the number of malicious agents or
access to labeled validation data.

Building on this core detection, SecFedRL integrates three supporting
components---reputation-aware aggregation, checkpoint-based collapse prevention, and
lightweight hash-chained consensus---forming a complete, fully decentralized
Byzantine-resilient FRL framework. The specific contributions are:
\begin{itemize}
  \item A fully decentralized Byzantine-resilient FRL framework tailored to the attack
        characteristics of reinforcement learning is proposed.
  \item A virtual-trajectory-based multi-metric consistency scoring mechanism is
        introduced for accurate malicious update detection without ground-truth labels.
  \item The effectiveness and generality of the proposed approach are validated through
        comprehensive reinforcement learning benchmarks.
\end{itemize}

%% ════════════════════════════════════════════════════════════════════════
%% §3  RELATED WORK
%% ════════════════════════════════════════════════════════════════════════
\section{Related Work}
\label{sec:related}

Extensive research has focused on enhancing the robustness, security, and reliability
of FL and FRL systems. This section reviews foundational FL frameworks, classical
Byzantine-resilient aggregation, attack modelling, cryptographic and
verification-based defences, communication-efficient strategies, consensus and
reputation mechanisms, and RL-specific approaches to Byzantine mitigation.

\subsection{Foundations of FL and Federated RL}

McMahan et al.~\cite{ref25} introduced the FedAvg algorithm for
communication-efficient collaborative training of deep networks from decentralized
data, establishing the foundational paradigm for FL. Yang et al.~\cite{ref24}
formalised the concept of federated machine learning and identified its core
application scenarios across horizontal, vertical, and transfer learning settings.
Kairouz et al.~\cite{ref2} review the progress and unresolved issues in FL, addressing
topics such as optimization, privacy, robustness, and fairness. All the aforementioned
papers elucidate the principal challenges in the domain of FL; however, none address
the particular attributes of Reinforcement Learning (RL), specifically the stochastic
nature of exploration and the non-IID distribution of trajectories. Zhu et
al.~\cite{ref3} examine the impact of data heterogeneity on the efficacy of FL,
asserting that data heterogeneity is the paramount concern, particularly in
reinforcement learning when environments produce trajectories.

\subsection{Classical Byzantine-Resilient Aggregation}

Initial studies on Byzantine-resilient distributed learning helped lay the theoretical
foundations for robust aggregation. Krum~\cite{ref19} was proposed as a
Byzantine-gradient descent rule that chooses update vectors near their neighbours.
Chen et al.~\cite{ref20} introduced a Byzantine-robust gradient descent algorithm
based on the geometric median-of-means that achieves a near-optimal estimation error
even when up to half of the machines are attacked. Xie et al.~\cite{ref21} expanded on
these approaches by introducing generalized Byzantine-tolerant SGD with better
robustness guarantees. Yin et al.~\cite{ref22} contributed to the theory of
Byzantine-robust learning by proving optimal statistical rates through coordinate-wise
median and trimmed-mean operations. Although useful, these approaches function purely
in the parameter space and require the assumption that benign gradients are
well-clustered---an assumption that does not hold in FRL owing to its highly variable
non-IID policy gradients.

\subsection{Attack Modelling and Robustness Evaluation}

Lyu et al.~\cite{ref23} provided a taxonomy of adversarial attacks to FL that
categorizes attacks as data poisoning, model poisoning, and inference attacks. Fang
et al.~\cite{ref16} demonstrated the existence of powerful local model-poisoning
attacks that can significantly degrade the performance of Byzantine-robust aggregation
rules such as Krum and trimmed mean. Shejwalkar and Houmansadr~\cite{ref8} provided an
optimal attack design for model poisoning attacks that maximizes attack efficacy while
minimizing detection likelihood. Shejwalkar et al.~\cite{ref5} revisited these attacks
to assess their feasibility in realistic FL environments. To counteract these attacks,
Chen et al.~\cite{ref27} designed FLRAM, which employs multiple aggregation filters to
protect against model-poisoning attacks.

\subsection{Cryptographic and Verification-Based Defences}

Bonawitz et al.~\cite{ref4} presented practical secure aggregation using
secret-sharing techniques, providing privacy protection while ensuring no single
contribution is revealed to the server. Cui et al.~\cite{ref26} devised a
dual-verifiable FL protocol based on vector commitment schemes, offering cryptographic
guarantees for verifying aggregation integrity from both server and client perspectives.
While these techniques provide strong theoretical security, the high computational and
communication costs make them impractical for FRL settings requiring near-real-time
policy updates.

\subsection{Communication-Efficient and Layer-Wise Aggregation}

Du et al.~\cite{ref12} presented high-dimensional stochastic gradient quantisation
algorithms for communication-efficient edge learning. Lan et al.~\cite{ref9} developed
an adaptive compression framework balancing learning accuracy and available computing
resources. Lin et al.~\cite{ref11} proposed ensemble distillation for robust model
aggregation during FL using knowledge distillation. Xu et al.~\cite{ref10} proposed a
layer-adaptive sparsified model aggregation algorithm achieving Byzantine-resilience.
Tegin and Duman~\cite{ref13} studied over-the-air aggregation in FL using
time-varying wireless channels. Despite mitigating communication overheads, most of
these approaches focus on supervised learning and do not consider sparse policy-level
adversarial attacks common in FRL.

\subsection{Consensus, Reputation, and Client Selection Mechanisms}

Sokhankhosh and Rouhani~\cite{ref6} present Proof-of-Collaborative-Learning, a
multi-winner consensus-based algorithm in a blockchain setting for validating FL
updates. Li and Lyu~\cite{ref7} prove convergence of sequential FL with heterogeneous
data, establishing theoretical fault-tolerance bounds. Deressa and Hasan~\cite{ref14}
propose TrustBandit, a multi-armed bandit framework for robust client selection. Ni
et al.~\cite{ref18} present rFedFW---a reputation-based aggregation method for
Byzantine-robust FL in IoT using a trust calculation based on historical contributions
and gradient consistency. These methods improve reliability and scalability, yet remain
predominantly focused on supervised FL with insufficient sensitivity to policy-space
behavior.

\subsection{Byzantine Resilience in Decentralized and Heterogeneous Settings}

Hou et al.~\cite{ref28} developed a Byzantine-robust decentralized FL framework which
evaluates feature extraction layer similarity and uses delayed aggregation, providing
robustness without a central coordinator. Xia et al.~\cite{ref29} demonstrate that
Byzantine-resilient methods can be developed for decentralized frameworks using
distributed optimization approaches such as PDMM. Karami et al.~\cite{ref1} developed
OptiGradTrust---a Byzantine-robust FL framework combining multi-feature gradient
analysis with RL-based trust assessment. However, all these techniques still fail to
address the sequential nature of FRL's decision-making.

\subsection{FRL-Specific Defences}

Fan et al.~\cite{ref15} provided the first theoretical treatment of fault-tolerant FRL,
establishing convergence guarantees under Byzantine failures. Xie and Song~\cite{ref17}
introduced FedKL, which penalises KL divergence during aggregation to address data
heterogeneity in federated RL. Salgia and Chi~\cite{ref30} analysed the fundamental
trade-off between sample complexity and communication cost in federated Q-learning.
Despite their contributions, these methods operate primarily in parameter space, lack
integrated mechanisms for collapse prevention, and do not perform behavioral Byzantine
detection.

\subsection{Summary and Research Gap}

Overall, the existing literature has made substantial progress in Byzantine-resilient
FL aggregation~\cite{ref19,ref20,ref21,ref22,ref27}, attack
modelling~\cite{ref5,ref8,ref16,ref23}, cryptographic
verification~\cite{ref4,ref26}, communication-efficient
strategies~\cite{ref9,ref10,ref11,ref12,ref13}, consensus and reputation
mechanisms~\cite{ref6,ref7,ref14,ref18}, decentralized
robustness~\cite{ref1,ref28,ref29}, and federated RL-specific
defences~\cite{ref15,ref17,ref30}. However, no existing approach provides a unified
framework combining (i) ground-truth-free policy-space behavioral Byzantine detection,
(ii) reputation-aware aggregation with decentralized consensus, (iii) proactive
catastrophic collapse prevention, and (iv) lightweight auditability---all within a
fully decentralized FRL architecture. SecFedRL fills these gaps with virtual-trajectory
consistency scoring, multi-metric behavioral analysis, reputation-based weighting,
checkpoint-based recovery, and hash-chained ledger consensus.

Table~\ref{tab:comparison} provides a comparative analysis of existing federated
learning defence mechanisms across seven key capability dimensions.

\begin{table}[htbp]
\centering
\caption{Comparison of Existing Approaches with SecFedRL}
\label{tab:comparison}
\footnotesize
\resizebox{\linewidth}{!}{%
\begin{tabular}{lcccccccc}
\toprule
\textbf{Approach} &
\makecell{\textbf{Byz.}\\\textbf{Robust}} &
\makecell{\textbf{Policy}\\\textbf{Space}} &
\makecell{\textbf{Rep.}\\\textbf{Aware}} &
\makecell{\textbf{Decen.}\\\textbf{Consns.}} &
\makecell{\textbf{Collapse}\\\textbf{Prevent.}} &
\makecell{\textbf{Audit.}\\\textbf{Ledger}} &
\makecell{\textbf{FRL}\\\textbf{Specific}} \\
\midrule
Krum~\cite{ref19}                   & \yes & \no  & \no  & \no  & \no  & \no  & \no  \\
Geo.\ Median~\cite{ref20}           & \yes & \no  & \no  & \no  & \no  & \no  & \no  \\
Trimmed Mean~\cite{ref22}           & \yes & \no  & \no  & \no  & \no  & \no  & \no  \\
Gen.\ Byz.\ SGD~\cite{ref21}        & \yes & \no  & \no  & \no  & \no  & \no  & \no  \\
FLRAM~\cite{ref27}                  & \yes & \no  & \no  & \no  & \no  & \no  & \no  \\
Fang et al.~\cite{ref16}            & \no  & \no  & \no  & \no  & \no  & \no  & \no  \\
Shejwalkar et al.~\cite{ref5,ref8}  & \yes & \no  & \no  & \no  & \no  & \no  & \no  \\
Bonawitz et al.~\cite{ref4}         & \no  & \no  & \no  & \no  & \no  & \yes & \no  \\
Cui et al.~\cite{ref26}             & \yes & \no  & \no  & \no  & \no  & \yes & \no  \\
Lan et al.~\cite{ref9}              & \no  & \no  & \no  & \no  & \no  & \no  & \no  \\
Xu et al.~\cite{ref10}              & \yes & \no  & \no  & \no  & \no  & \no  & \no  \\
Sokhankhosh et al.~\cite{ref6}      & \no  & \no  & \no  & \yes & \no  & \yes & \no  \\
TrustBandit~\cite{ref14}            & \yes & \no  & \yes & \no  & \no  & \no  & \no  \\
rFedFW~\cite{ref18}                 & \yes & \no  & \yes & \no  & \no  & \no  & \no  \\
Li \& Lyu~\cite{ref7}               & \no  & \no  & \no  & \no  & \no  & \no  & \no  \\
OptiGradTrust~\cite{ref1}           & \yes & \no  & \yes & \no  & \no  & \no  & \no  \\
Hou et al.~\cite{ref28}             & \yes & \no  & \no  & \yes & \no  & \no  & \no  \\
Xia et al.~\cite{ref29}             & \yes & \no  & \no  & \no  & \no  & \no  & \no  \\
Fan et al.~\cite{ref15}             & \yes & \no  & \no  & \no  & \no  & \no  & \yes \\
FedKL~\cite{ref17}                  & \no  & \no  & \no  & \no  & \no  & \no  & \yes \\
Salgia \& Chi~\cite{ref30}          & \no  & \no  & \no  & \no  & \no  & \no  & \yes \\
\midrule
\textbf{SecFedRL (Ours)}            & \yes & \yes & \yes & \yes & \yes & \yes & \yes \\
\bottomrule
\end{tabular}%
}
\end{table}

\smallskip\noindent\textit{Note: The evaluation dimensions in Table~\ref{tab:comparison}
were defined by the authors to characterize the specific gaps that SecFedRL addresses.
Comparisons reflect design choices and architectural properties rather than independent
third-party benchmarking. Readers should consult the cited works directly when assessing
the full capability scope of each baseline method.}

%% ════════════════════════════════════════════════════════════════════════
%% §4  METHODOLOGY
%% ════════════════════════════════════════════════════════════════════════
\section{Methodology}
\label{sec:methodology}

\subsection*{The SecFedRL Framework}

To address the limitations identified in prior work, this section presents SecFedRL, a
Byzantine-resilient framework for P2P FRL. SecFedRL enables strong collaborative
policy learning in adversarial environments through a fully decentralized architecture.
It integrates behavioral Byzantine detection, reputation-aware aggregation, catastrophic
collapse prevention, and performance-based consensus, all without needing central
coordination or ground-truth supervision. The overall system architecture is shown in
Figure~\ref{fig:architecture}. The training algorithm is summarized in
Algorithm~\ref{alg:secfedrl}, and the hyperparameters are listed in
Table~\ref{tab:hyperparams}.

\subsection{Threat Model}

SecFedRL is designed assuming that Byzantine adversaries can control local model
updates, manipulate reward structures or gradients in arbitrary ways, and inflict
malicious attacks (by scaling updates, sign-flipping updates, conducting
policy-targeted poisonings, and/or coordinated attacks) on local agents. Byzantine
adversaries can control no more than 40\% of participating agents ($f < 0.4n$), which
is consistent with all currently known practical threat models in distributed RL.
Additionally, Byzantine malicious adversaries generally do not have access to
cryptographic mechanisms (such as SHA-256 hashing) and cannot maintain control of the
consensus mechanism ($2n/3$ threshold). Neither a trusted server nor ground-truth
labels need to be trusted; this yields an honest/respected majority assumption
($f < n/2$) for all agents. Byzantine adversaries will have only grey-box information
about the agents' learning protocol and will not know their own private trajectories or
internal reputation scores. The goal of manipulating or compromising other agent
mechanisms is to lower the overall performance of the collective policy or create
policy instability while training through coordinated or purely-infliction-based means.

\noindent\textbf{Environment Homogeneity Assumption.}
SecFedRL's virtual trajectory--based behavioral analysis (Section~\ref{sec:detection})
requires that participating agents operate within the same or structurally identical
environment instance, ensuring that states sampled from agent $i$'s replay buffer
constitute valid behavioral evaluation inputs for peer $j$'s policy. Under this
assumption, any divergence in peer $j$'s action distribution over shared states can be
attributed to policy differences rather than environment heterogeneity. This assumption
holds for the CartPole-v1, LunarLander-v2, and Pendulum-v1 evaluations reported here,
where all agents interact with the same gym environment class. Extending SecFedRL to
genuinely heterogeneous environments (e.g., multi-task settings or sim-to-real
transfer) would require environment-invariant behavioral representations and is
identified as a direction for future work (Section~\ref{sec:conclusion}).

\subsection{System Overview}

The decentralized model of SecFedRL for all entities runs as a P2P mesh topology where
each agent $i \in [1,N]$ communicates with $k$ randomly selected peers
($k \in [2,6]$), resulting in an $O(kn)$ communication complexity that is independent
of the total number of agents $N$. Each training round $t$ includes seven phases:
\begin{enumerate}
  \item Local agent A2C training on agent-specific trajectories
  \item Agent-to-agent P2P model exchange with $k$ peers
  \item Behavioral analysis of virtual trajectory data to identify peer behaviors
  \item Reputation-aware aggregation of validated and approved peer updates
  \item Collapse detection and local checkpoint recovery
  \item Performance-based consensus voting
  \item Hash-chained ledger recording
\end{enumerate}
The entire working flow is explained in Algorithm~\ref{alg:secfedrl}.

\begin{figure}[htbp]
  \centering
  \includegraphics[width=0.95\linewidth]{image_04.png}
  \caption{System Architecture for SecFedRL Framework.}
  \label{fig:architecture}
\end{figure}

Figure~\ref{fig:architecture} depicts the overall architecture of the proposed
SecFedRL system which allows secure distributed FRL operation by decentralized agents
in the presence of Byzantine adversaries. The framework consists of multiple
distributed autonomous clients (agents) that collaboratively develop an RL model
through P2P communication, without depending on a central aggregation server. On the
client side, each agent has a local A2C model comprised of an actor and a critic
network, and the agents perform local training by interacting with their respective
environments and creating trajectories. After local training, clients perform a P2P
exchange of RL model updates with neighbouring agents using an assessment methodology
based on policy-space behavioral analysis, rather than simply comparing model
parameters. The system uses a multi-metric Byzantine detection mechanism that evaluates
behavioral consistency among peer models using seven metrics: KL divergence, cosine
similarity, overall policy agreement, value alignment, entropy divergence, reverse KL
divergence, and gradient magnitude consistency.

The behavioral evaluation filters out suspicious updates, with remaining trustworthy
updates passed to the reputation-aware aggregation module. Each client maintains
dynamic reputation scores assigned to each peer based on historical performance,
behavioral consistency, and gradient alignment. The validated updates are combined
using a weighted average to produce an aggregated model. The framework continuously
monitors performance to detect sudden drops in average reward. If detected, the system
recovers the model from previously saved checkpoints. Finally, verification of
aggregated model updates is accomplished through a hash-chained audit ledger that records
each accepted round as a SHA-256 hash of the aggregated model state, chained to the
previous accepted round's hash to provide transparency and tamper-evident audit trails.

\begin{figure}[htbp]
  \centering
  \includegraphics[width=0.95\linewidth]{image_02.png}
  \caption{Round-wise execution time.}
  \label{fig:roundtime}
\end{figure}

In Figure~\ref{fig:roundtime}, per-round execution time increases gradually as training
progresses, from approximately 25--50~s in early rounds (when policies are random and
behavioral analysis is fast) to 140--175~s in later rounds (when learned policies are
more complex and the 7-metric analysis requires fuller computation). The 10-round
moving average confirms this upward trend, stabilising at approximately 155--170~s by
round~40. The distribution (right panel) shows a mean of 139.33~s and median of
143.35~s, with a range of 25--225~s, capturing round-to-round variance due to
non-synchronous P2P communication. This per-round overhead is consistent with the
4.2--8.3\% overhead reported in Table~\ref{tab:scalability}, confirming that the
behavioral detection plugin imposes a bounded and proportional cost within acceptable
limits for the federated training loop.

\subsection{Local Actor-Critic Training}

A2C with GAE is selected for its established use in FRL convergence analysis~\cite{ref15}
and its tractable gradient variance under the trajectory heterogeneity characteristic of
P2P FRL. Policy-gradient variance is a primary convergence obstacle in decentralized
settings; GAE's $\lambda$-weighted advantage ($\lambda_{\text{GAE}}=0.95$, $\gamma=0.99$)
provides a principled bias-variance trade-off that stabilizes updates across agents with
different local data distributions. Alternative algorithms such as PPO and SAC---while
powerful in single-agent and centralized settings---introduce additional hyperparameters
(clipping ratio, temperature) whose tuning is poorly conditioned under the
Byzantine-induced distribution shift that SecFedRL must handle. A2C's simpler update rule
facilitates reproducible behavioral profiling across agents, which is essential for the
virtual trajectory consistency scoring in Section~\ref{sec:detection}.

Each agent trains its own A2C model based on Generalized Advantage Estimation (GAE)
using its own independent experience. The agents' model policy $\pi(a|s)$ and value
function $V(s)$ share a convolutional feature extraction network with separate actor
and critic output networks. The optimization objective is:
\begin{equation}
  L(\theta) = L_\pi(\theta) + c_v \cdot L_V(\theta) - c_e \cdot H(\pi(\cdot|s))
  \label{eq:loss}
\end{equation}
where:
\begin{itemize}
  \item $L_\pi(\theta)$ = policy gradient loss = $-\mathbb{E}_t\bigl[\log\pi(a_t|s_t;\theta)\cdot\hat{A}_t\bigr]$, where $\hat{A}_t$ is the generalized advantage estimate
  \item $L_V(\theta)$ = value function loss = $\mathbb{E}_t\bigl[(V(s_t;\theta) - V^{\text{target}}_t)^2\bigr]$
  \item $H(\pi(\cdot|s))$ = policy entropy regularization = $-\sum_a \pi(a|s)\log\pi(a|s)$
  \item $c_v$ = critic loss weight (value: 0.5)
  \item $c_e$ = entropy coefficient (value: 0.01, see Table~\ref{tab:hyperparams})
\end{itemize}

\begin{figure}[htbp]
  \centering
  \includegraphics[width=0.95\linewidth]{image_01.png}
  \caption{Training metrics of SecFedRL across various rounds.}
  \label{fig:training}
\end{figure}

Figure~\ref{fig:training} shows that the cumulative rewards of agents demonstrate a
continuous upward trend. As agents learn and implement optimal policies over time, the
reward trend is steadily increasing. The episode step trend illustrates that with each
increment in rounds, agents' episode steps are converging to be stable. Variability
among agents in different environments results in temporary variances in performance.
However, as the reward standard deviation progressively stabilises, this indicates that
agents' policies update consistently, supporting the stability and convergence of the
local A2C training method.

\subsection{Multi-Metric Behavioral Byzantine Detection}
\label{sec:detection}

SecFedRL uses a plugin-based behavioral test suite to evaluate peer models on 100
procedurally generated states, $s' \sim \mathcal{S}_{\text{test}}$, drawn from the
agent's replay buffer. Seven complementary metrics capture policy divergence,
behavioral inconsistency, critic corruption, and parameter anomalies:

\noindent\textbf{KL divergence:}
\[
  D_{\text{KL}}(\pi_i \| \pi_j)(s') = \sum_a \pi_i(a|s') \cdot \log\frac{\pi_i(a|s')}{\pi_j(a|s')}
\]

\noindent\textbf{Cosine similarity:}
\[
  \cos(\Phi_i, \Phi_j) = \frac{\Phi_i \cdot \Phi_j}{\|\Phi_i\| \cdot \|\Phi_j\|},\quad
  \text{where } \Phi_i = f_{\theta_i}(s')
\]

\noindent\textbf{Policy agreement:}
\[
  P_{\text{agree}}(\pi_i, \pi_j) =
    \mathbb{E}_{s'}\!\left[\mathbf{1}\!\left(\argmax_a \pi_i(a|s') = \argmax_a \pi_j(a|s')\right)\right]
\]

\noindent\textbf{Value alignment:}
\[
  V_{\text{align}}(i,j) =
    \max\!\left(0,\; 1 - \frac{|V_i(s') - V_j(s')|}{\max(|V_i(s')|,\, |V_j(s')|) + \varepsilon}\right)
\]

\noindent\textbf{Entropy divergence:}
\[
  E_{\text{div}}(i,j) = \bigl|H(\pi_i(\cdot|s')) - H(\pi_j(\cdot|s'))\bigr|,\quad
  \text{where } H(\pi) = -\sum_a \pi(a|s')\log\pi(a|s')
\]

\noindent\textbf{Reverse KL:}
\[
  D_{\text{KL}}(\pi_j \| \pi_i)(s') = \sum_a \pi_j(a|s') \cdot \log\frac{\pi_j(a|s')}{\pi_i(a|s')}
\]

\noindent\textbf{Gradient magnitude consistency:}
\[
  G_{\text{mag}}(i,j) = \bigl|\|\nabla_{\theta_i} L(\theta_i)\|_2 - \|\nabla_{\theta_j} L(\theta_j)\|_2\bigr|
\]

The composite suspicion score for peer $j$ is defined as:
\begin{equation}
  S_j = \sum_{m=1}^{7} w_m \cdot \frac{M_m(j) - \mu_m}{\sigma_m}
  \label{eq:suspicion}
\end{equation}
where:
\begin{itemize}
  \item $S_j$ = final suspicion score for peer $j$ (higher = more suspicious)
  \item $M_m(j)$ = value of the $m$-th behavioral metric computed for peer $j$
  \item $\mu_m$ = mean value of metric $m$ across all participating peers in round $t$
  \item $\sigma_m$ = standard deviation of metric $m$ across all participating peers
  \item $w_m$ = weight assigned to the $m$-th metric ($\sum w_m = 1$; equal weights $w_m = 1/7$ by default)
  \item $m \in \{\text{KL-div, Cosine, Policy-agree, Value-align, Entropy-div, Rev-KL, Grad-mag}\}$
\end{itemize}

The composite suspicion score combines multiple behavioral signals to identify unusual
peer updates. The weighted combination helps SecFedRL find malicious updates even when
attackers try to manipulate only certain parameters or gradients.

\noindent\textbf{Statistical stabilization via sliding window.}
With $k=4$ peers per agent, the per-round sample estimates of $\mu_m$ and $\sigma_m$
are computed over a sliding window of the last 10 rounds, accumulating up to
$10 \times k = 40$ observations per metric per agent. This window size was chosen to
balance responsiveness to newly injected Byzantine peers against estimate stability:
fewer than 10 rounds yields high-variance $\sigma_m$ estimates that inflate false
positive rates during the warm-up phase, while more than 10 rounds delays detection
of late-joining Byzantine agents. The first 10 rounds use all available observations;
from round 11 onward, only the most recent 10 rounds contribute to $\mu_m$ and $\sigma_m$.

\subsection{Reputation-Aware Aggregation}

SecFedRL maintains ground-truth-free reputation scores $r_j(t) \in [0,1]$ for each
peer $j$, updated each round using:
\begin{equation}
  r_j(t) = \alpha \cdot r_j(t-1) + (1-\alpha)\cdot
    \bigl[\beta_1 \cdot \hat{R}_j + \beta_2 \cdot C_j + \beta_3 \cdot G_j\bigr]
  \label{eq:reputation}
\end{equation}
where:
\begin{itemize}
  \item $\hat{R}_j$ = normalized episodic reward of peer $j$ in round $t$ (clipped to $[0,1]$)
  \item $C_j$ = behavioral consistency score ($1 - S_j / \max_k S_k$)
  \item $G_j$ = gradient alignment score (cosine similarity of $\nabla\theta_j$ and $\nabla\theta_g$)
  \item $\alpha = 0.9$ = exponential moving average decay factor
  \item $\beta_1 = \beta_2 = \beta_3 = 1/3$ = component weights (equal weighting)
  \item Initial reputation $r_j(0) = 0.5$ for all peers
\end{itemize}

The final aggregation weights combine short-term performance with long-term peer
reliability:
\begin{equation}
  w_j = \frac{r_j(t)\cdot(1 - S_j)}{\displaystyle\sum_{k \in \mathcal{B}_t}
    \bigl[r_k(t)\cdot(1 - S_k)\bigr]}
  \label{eq:weights}
\end{equation}
where $\mathcal{B}_t$ is the trusted benign set after Byzantine filtering (Phase 4,
Algorithm~\ref{alg:secfedrl}). The global model is then updated according to:
\begin{equation}
  \theta_g = \sum_{j \in \mathcal{B}_t} w_j \cdot \theta_j
  \label{eq:global}
\end{equation}
where the weights $w_j$ already sum to 1 by construction (Eq.~\ref{eq:weights}),
so no additional normalization is required.

\subsection{Catastrophic Policy Collapse Prevention}

To avoid permanent policy corruption from poisoned aggregations, SecFedRL uses
checkpoint-based recovery. The recovery process merges the current policy with a
previously stored checkpoint:
\begin{equation}
  \theta_{\text{recover}} = \lambda \cdot \theta_{\text{checkpoint}}
    + (1-\lambda)\cdot\theta_{\text{current}}
  \label{eq:recovery}
\end{equation}
where $\lambda = \min(\lambda_{\max},\, \Delta r_t / \bar{r}_{t-1})$ scales
proportionally with the normalized reward drop severity
($\Delta r_t = \bar{r}_{t-1} - \bar{r}_t$), capped at $\lambda_{\max} = 0.85$
(see Table~\ref{tab:hyperparams}) to prevent full checkpoint replacement. A collapse
is defined as $\Delta r_t > 20\%$ of the prior moving average. Checkpoints are saved
every 5 rounds. Recovery takes place without interrupting training.

\subsection{Performance-Based Consensus and Ledger}

Auditability is ensured through a simple hash-chained ledger. After aggregation,
agents vote pass or fail based on a \textbf{relative performance improvement
threshold}. Rather than using a fixed absolute reward value (which would not generalize
across environments with different reward scales, such as the negative-valued rewards
in Pendulum-v1), the consensus criterion is defined as:
\begin{equation}
  D_t = \begin{cases}
    1 & \text{if } \displaystyle\sum_i \mathrm{vote}_i(t)
        \geq \left\lceil\tfrac{2}{3}\cdot N\right\rceil \\[4pt]
    0 & \text{otherwise}
  \end{cases}
  \label{eq:consensus}
\end{equation}
where $\mathrm{vote}_i(t) = 1$ if agent $i$ observes that the aggregated reward for
round $t$ improves by at least $\delta = 5\%$ over the 5-round moving average
$\bar{r}_{t-5:t-1}$, and $\mathrm{vote}_i(t) = 0$ otherwise. This relative
formulation is environment-agnostic: for CartPole-v1 (rewards $\approx 400$--500), the
threshold tracks absolute improvement; for Pendulum-v1 (rewards $\approx -185$ to
$-847$), it tracks percentage reduction in the magnitude of the negative reward.

\noindent Where:
\begin{itemize}
  \item $D_t$ = final consensus decision for round $t$ (1 = accepted, 0 = rejected/rolled back)
  \item $\mathrm{vote}_i(t)$ = vote from agent $i$ (relative improvement criterion above)
  \item $N$ = total number of participating agents
\end{itemize}

The 5\% improvement margin $\delta$ was chosen empirically to balance sensitivity to
genuine progress against noise from individual-round variance (see Table~\ref{tab:hyperparams}
for the fixed hyperparameter used in the consensus vote count threshold).

A training round is accepted only if at least two-thirds of the participating agents
approve the aggregated update, ensuring Byzantine tolerance and decentralized
decision-making. Accepted rounds are recorded as immutable ledger entries with
${<}0.1\%$ overhead.

\subsection{SecFedRL Training Algorithm}

\begin{algorithm}[htbp]
\caption{SecFedRL Full Round-Level Training}
\label{alg:secfedrl}
\begin{algorithmic}[1]
\Require Local policy parameters $\{\theta_i\}$ for $i = 1$ to $N$; global model $\theta_g$
\Require $N_{\text{vt}} = 100$ (virtual trajectory states per behavioral evaluation)
\Require $\delta_S$: suspicion threshold (dynamic: initialized to $2\sigma$ above peer mean, updated each round)
\Require Peers per agent $k = 4$; relative consensus margin $\delta = 5\%$; checkpoint interval $C_{\text{int}} = 5$ rounds
\Ensure Updated global model $\theta_g$ robust to Byzantine updates
\State \textbf{Note:} SecFedRL requires no prior knowledge of the number of malicious agents; Byzantine/benign separation uses dynamic statistical thresholding.
\For{each round $t = 1$ to $T$}
  \State \textbf{--- Phase 1: Local A2C Training ---}
  \State Each agent $i$ trains its local A2C model on its own trajectories using $L = L_\pi + c_v\cdot L_V - c_e\cdot H(\pi)$ (Eq.~\ref{eq:loss}).
  \State \textbf{--- Phase 2: P2P Model Exchange ---}
  \State Agent $i$ randomly selects $k$ peers and exchanges local policy weights $\theta_i$ with each selected peer $j$.
  \State \textbf{--- Phase 3: Virtual Trajectory Behavioral Analysis ---}
  \For{each peer $j$ in the received peer set}
    \State Sample $N_{\text{vt}}$ states $\{s'_1,\ldots,s'_{N_{\text{vt}}}\}$ from agent $i$'s replay buffer.
    \For{each state $s'_{\text{vt}}$}
      \State Compute: $\Phi_i(s'_{\text{vt}}) = f_{\theta_i}(s'_{\text{vt}})$ and $\Phi_j(s'_{\text{vt}}) = f_{\theta_j}(s'_{\text{vt}})$
    \EndFor
    \State Compute all 7 behavioral consistency metrics for peer $j$ (Section~\ref{sec:detection}):
    \State \hspace{1em}$D_{\text{KL}}(\pi_i\|\pi_j)$, $\cos(\Phi_i,\Phi_j)$, $P_{\text{agree}}$, $V_{\text{align}}$, $E_{\text{div}}$, $D_{\text{KL}}(\pi_j\|\pi_i)$, $G_{\text{mag}}$
    \State Compute composite suspicion score $S_j = \sum_m w_m\cdot(M_m(j)-\mu_m)/\sigma_m$ (Eq.~\ref{eq:suspicion})
  \EndFor
  \State \textbf{--- Phase 4: Byzantine Filtering and Reputation-Aware Aggregation ---}
  \State Exclude peer $j$ if $S_j > \mu_S + \delta_S$; collect trusted benign set $\mathcal{B}_t$.
  \For{each $j \in \mathcal{B}_t$}
    \State Update reputation: $r_j(t) = \alpha\cdot r_j(t{-}1) + (1{-}\alpha)\cdot[\beta_1\hat{R}_j + \beta_2 C_j + \beta_3 G_j]$ (Eq.~\ref{eq:reputation})
  \EndFor
  \State Compute weights: $w_j = r_j(t)\cdot(1-S_j)\;/\;\sum_{k\in\mathcal{B}_t}[r_k(t)\cdot(1-S_k)]$ (Eq.~\ref{eq:weights})
  \State Update global model: $\theta_g \leftarrow \sum_{j\in\mathcal{B}_t} w_j\,\theta_j$ (Eq.~\ref{eq:global})
  \State \textbf{--- Phase 5: Catastrophic Collapse Detection and Recovery ---}
  \If{$(\bar{r}_{t-1} - \bar{r}_t) > \text{collapse\_threshold}$}
    \State $\theta_g \leftarrow \lambda\cdot\theta_{\text{checkpoint}} + (1-\lambda)\cdot\theta_g$,\; $\lambda = \min(\lambda_{\max},\,\Delta r_t / \bar{r}_{t-1})$, $\lambda_{\max}=0.85$ (Eq.~\ref{eq:recovery})
  \EndIf
  \State Save checkpoint every $C_{\text{int}} = 5$ rounds.
  \State \textbf{--- Phase 6: Performance-Based Consensus Vote ---}
  \State Each agent $i$: $\mathrm{vote}_i(t) = 1$ if round reward improves by $\geq\delta = 5\%$ over prior 5-round moving average, else 0.
  \State Consensus (Eq.~\ref{eq:consensus}): $D_t = 1$ if $\sum_i \mathrm{vote}_i(t) \geq \lceil\tfrac{2}{3}N\rceil$, else roll back $\theta_g$.
  \State \textbf{--- Phase 7: Hash-Chained Ledger Recording ---}
  \If{$D_t = 1$}
    \State Append $\{t,\;\mathrm{hash}(\theta_g),\;\sum_i \mathrm{vote}_i(t),\;|\mathcal{B}_t|\}$ to the immutable ledger.
  \EndIf
\EndFor
\State \Return $\theta_g$
\end{algorithmic}
\end{algorithm}

\begin{table}[htbp]
\centering
\caption{Hyperparameters.}
\label{tab:hyperparams}
\begin{tabular}{lll}
\toprule
\textbf{Parameter} & \textbf{Value} & \textbf{Description} \\
\midrule
$N$              & 20                 & Total number of agents \\
$k$              & 4                  & Peers per agent per round \\
$\gamma$         & 0.99               & Discount factor (A2C/GAE) \\
$\lambda_{\text{GAE}}$ & 0.95         & GAE trace-decay parameter (advantage estimation) \\
$c_v$            & 0.5                & Critic loss weight in A2C loss (Eq.~\ref{eq:loss}) \\
$c_e$            & 0.01               & Entropy coefficient in A2C loss (Eq.~\ref{eq:loss}) \\
$\eta$           & $3\times10^{-4}$   & Learning rate (Adam optimizer) \\
$N_{\text{vt}}$  & 100                & Virtual trajectory states per behavioral evaluation \\
$\lambda_{\max}$ & 0.85               & Maximum checkpoint recovery interpolation weight (Eq.~\ref{eq:recovery}) \\
$\delta$         & 5\%                & Relative reward improvement for consensus vote pass (environment-agnostic; replaces fixed $\tau=300$; Eq.~\ref{eq:consensus}) \\
$T_{\text{local}}$ & 128--512 steps & Local training steps per communication round per agent (environment-dependent; see \url{https://github.com/Yog37h/byzsecure}) \\
\bottomrule
\end{tabular}
\end{table}

%% ════════════════════════════════════════════════════════════════════════
%% §5  RESULTS AND DISCUSSION
%% ════════════════════════════════════════════════════════════════════════
\section{Results and Discussion}
\label{sec:results}

This section presents a comprehensive evaluation of SecFedRL within the
context of adversarial FRL. Several features were included in the assessment including
robustness to Byzantine attacks, detection accuracy, adaptability to various
environments, scalability, and capability of contributing components. The evaluation
was performed using the set of environments from OpenAI Gym to incorporate multiple
characteristics of reinforcement learning such as discrete vs.\ continuous, sparse
vs.\ dense rewards, and control tasks.

\subsection{Experimental Setup}

The learning algorithm, SecFedRL, was tested on three separate environments:
CartPole-v1, LunarLander-v2, and Pendulum-v1 using twenty independent agents with
20\% Byzantine fraction (unless otherwise specified). Each agent used A2C with GAE.
The overall experiment was conducted over 100 communication rounds, replicated across
five different stochastic samples. The Byzantine agents intentionally executed
Untargeted Model Poisoning (UMP) and Targeted Model Poisoning (TMP) attacks. Baseline
learning algorithms included FedAvg (no defense), Krum, and Trimmed Mean. Hyperparameter
descriptions are found in Table~\ref{tab:hyperparams}.

\textbf{Baseline Selection Rationale.} FedAvg~\cite{ref25}, Krum~\cite{ref19}, and
Trimmed Mean~\cite{ref22} were selected as baselines because they represent the
foundational layer of Byzantine-resilient aggregation and are the most widely compared
methods in the FRL robustness literature~\cite{ref15,ref5}. Two additional candidates
were reviewed but excluded: FedKL~\cite{ref17} addresses data \emph{heterogeneity} in
FRL via KL-divergence regularization of local policy updates; it is not a Byzantine
defence mechanism and contains no adversary detection component, making it an
inappropriate robustness baseline. rFedFW~\cite{ref18} is designed for IoT scenarios
with a trusted central server, which is architecturally incompatible with SecFedRL's
fully decentralized P2P topology. OptiGradTrust~\cite{ref1} was published as an arXiv
preprint in July 2025, after the experimental phase of this work was completed,
precluding direct experimental comparison; additionally, its design requires
centralized gradient collection from a trusted aggregation server, architecturally
incompatible with SecFedRL's serverless P2P mesh. FLRAM~\cite{ref27} applies
parameter-space median-based filtering that is mechanistically comparable to Krum and
Trimmed Mean; its performance characteristics relative to those two baselines are
characterized in~\cite{ref27} itself, so including it would not add an independent
evaluation dimension. Comparison against OptiGradTrust, FLRAM, and other recently
published defences is an important direction for future work.

\subsection{Overall Performance Under Adversarial Conditions}

The first experiment tests the learning performance of SecFedRL against 20\% Byzantine
agents using UMP attacks in the CartPole-v1 environment. The results are presented in
Table~\ref{tab:cartpole_perf}.

\begin{table}[htbp]
\centering
\caption{Performance Comparison Under 20\% Byzantine Attacks (CartPole-v1).}
\label{tab:cartpole_perf}
\begin{tabular}{lcccc}
\toprule
\textbf{Method} & \textbf{Final Reward} & \textbf{Success Rate} &
\textbf{Convergence Round} & \textbf{Overhead} \\
\midrule
FedAvg (No Defense)  & $78.0 \pm 12.3$    & 0\%  & Failed & 0\%   \\
Krum                 & $312.5 \pm 45.7$   & 62\% & $>$50  & 8.2\% \\
Trimmed-Mean         & $287.3 \pm 52.1$   & 57\% & $>$50  & 6.1\% \\
SecFedRL (Ours)      & $485.13 \pm 8.2$   & 97\% & 5      & 4.2\% \\
\bottomrule
\end{tabular}
\end{table}

\begin{figure}[htbp]
  \centering
  \includegraphics[width=0.95\linewidth]{image_07.png}
  \caption{The Convergence of episodic rewards under 20\% Byzantine attacks.}
  \label{fig:convergence}
\end{figure}

Figure~\ref{fig:convergence} shows SecFedRL's convergence speed analysis: panel (a)
plots the number of rounds required to reach each reward threshold, confirming that
reward thresholds up to 400 are achieved within 1 round while reaching threshold 450
requires \textbf{5 rounds} (``Final: 5 rounds'' annotation). Panel (b) shows the
corresponding time efficiency ($\approx$7.8 minutes at the 450 threshold). Panel (c)
shows steady learning efficiency of 15 reward/round during the convergent phase.
Panel (d) shows resource utilisation normalised to the highest threshold.

The 521\% improvement in final episodic reward over the FedAvg baseline is computed as:
\begin{equation}
  \text{Improvement\%} = \frac{R_{\text{SecFedRL}} - R_{\text{FedAvg}}}{|R_{\text{FedAvg}}|}
    \times 100\%
  \label{eq:improvement}
\end{equation}
where $R_{\text{SecFedRL}}$ and $R_{\text{FedAvg}}$ indicate the final episodic rewards
for the respective methods. According to these results, SecFedRL shows a 521\%
improvement in final episodic reward over the undefended FedAvg baseline, achieving
485.13 out of a maximum possible 500. SecFedRL achieved convergence (defined as first round where
the 5-round moving average reward exceeds the CartPole solved threshold of 195) within
5 rounds, as confirmed by Figure~\ref{fig:convergence}'s convergence-speed analysis
(threshold 450 annotation: ``Final: 5 rounds''), whereas parameter-space aggregation
techniques took longer than 50 rounds. Policy-loss components stabilize further by
round 25+, as shown in Figure~\ref{fig:loss}. The fast reward convergence of SecFedRL
is due to the behavioral filtering mechanism removing poisoned updates from early rounds,
allowing clean aggregation of locally pre-trained policies and preventing gradient
corruption that delays convergence in undefended or parameter-space-only methods.

\subsection{Byzantine Detection Accuracy}

All methods implemented for detection of Byzantine behavior have consistently produced
results showing successful Byzantine attack identification regardless of the environment
or type of attack. Table~\ref{tab:detection} reports the cumulative detection
performance metrics---defined as the fraction of Byzantine agents identified at least
once within the stated detection-round window. Figure~\ref{fig:detection} provides the
round-level temporal view.

\begin{table}[htbp]
\centering
\caption{Byzantine Detection Performance: Cumulative Detection (20\% Attack Fraction). TPR and FPR reflect cumulative detection---each Byzantine agent flagged at least once within the stated detection-round window. For average per-round detection rates, see Table~\ref{tab:byzantine_frac}.}
\label{tab:detection}
\begin{tabular}{llcccccl}
\toprule
\textbf{Environment} & \textbf{Attack} & \textbf{Cumul.\ TPR} & \textbf{Cumul.\ FPR} &
\textbf{Precision} & \textbf{Recall} & \textbf{F1} & \textbf{Detection Round} \\
\midrule
CartPole    & UMP & 1.00 & 0.00 & 1.00 & 1.00 & 1.00 & 25--35 \\
CartPole    & TMP & 1.00 & 0.00 & 1.00 & 1.00 & 1.00 & 28--38 \\
LunarLander & UMP & 1.00 & 0.00 & 1.00 & 1.00 & 1.00 & 30--42 \\
LunarLander & TMP & 1.00 & 0.00 & 1.00 & 1.00 & 1.00 & 32--45 \\
Pendulum    & UMP & 1.00 & 0.00 & 1.00 & 1.00 & 1.00 & 22--35 \\
Pendulum    & TMP & 1.00 & 0.00 & 1.00 & 1.00 & 1.00 & 25--38 \\
\bottomrule
\end{tabular}
\end{table}

\smallskip\noindent\textit{\textbf{Scope note:} Table~\ref{tab:detection} reports
cumulative detection performance over all 100 rounds (each Byzantine agent flagged at
least once by the specified detection round). Figure~\ref{fig:detection} shows the
round-level temporal dynamics using a 10-agent configuration (20\%\,$\times$\,10 = 2
Byzantine clients) selected for visual clarity; the 20-agent main experiments (4
Byzantines) show qualitatively identical phase structure (warm-up $\to$ variable $\to$
stabilised) and are summarised via Table~\ref{tab:detection}. Detection is variable
during the initial warm-up phase, then stabilises after the detection round specified
in the table. The TPR=1.00 in Table~\ref{tab:detection} reflects that all Byzantine
agents are eventually identified within the stated detection-round window; it does not
imply every individual round achieves 100\% detection. Figure~\ref{fig:detection}'s
statistics (Average: 1.5; Zero Detection: 4/50 rounds (8\%); Perfect Detection: 30/50
rounds (60\%)) confirm that the majority of rounds achieve full 2-of-2 detection, with
brief lapses only in the warm-up phase. These conditions represent the moderate-Byzantine
regime. As shown in Section~\ref{sec:byzantine_frac} (Table~\ref{tab:byzantine_frac}),
detection performance degrades at higher fractions: TPR falls to 0.95 at 40\% Byzantine
fraction under UMP. ROC analysis under varying suspicion threshold values is a natural
extension and is left as future work. A specific known failure mode is \textbf{mimicry
attack}: an adversary that deliberately calibrates its policy updates to remain within
the benign behavioral distribution could reduce the per-round detection rate below the
observed averages reported here. Under the UMP/TMP attack models tested, adversaries do
not perform such calibration, so the warm-up-phase lapses visible in
Figure~\ref{fig:detection} are transient and self-correcting. Resistance to adversarial
mimicry attacks represents an open limitation of the current detection framework and is
a priority direction for future work.}

\begin{figure}[htbp]
  \centering
  \includegraphics[width=0.95\linewidth]{image_06.png}
  \caption{Byzantine client detection across communication rounds.}
  \label{fig:detection}
\end{figure}

The combined reputation-and-suspicion framework distinguishes Byzantine agents
(reputation scores 0.42--0.73) from benign agents (reputation scores 0.64--0.89).
While raw reputation scores partially overlap in the 0.64--0.73 boundary region, the
composite multi-metric suspicion score $S_j$ provides a separating margin exceeding
0.15 standard deviations above the benign peer mean, enabling correct classification
in all test cases without ground-truth labels. A warm-up period of approximately
25--30 rounds was sufficient to establish stable behavioral baselines, after which the
framework enters its Late Training phase (round 30+) with consistent 2-of-2 Byzantine
detection. As shown in Figure~\ref{fig:detection}, detection is variable during the
warm-up and early training phases before stabilising after round 25--35 (CartPole-v1),
at which point all Byzantine agents have been cumulatively identified. The resulting
cumulative TPR=1.00 (Table~\ref{tab:detection}) reflects that no Byzantine agent evades
detection across the full 100-round training window.

\subsection{Cross-Environment Generalization}

To assess generalization capability, SecFedRL was evaluated across environments with
varying dynamics and reward structures under 20\% UMP attacks.
Table~\ref{tab:multienv} summarizes the results.

\begin{table}[htbp]
\centering
\caption{Multi-Environment Performance.}
\label{tab:multienv}
\begin{tabular}{lcccc}
\toprule
\textbf{Environment} & \textbf{Baseline Reward} & \textbf{SecFedRL Reward} &
\textbf{Improvement} & \textbf{Task Success} \\
\midrule
CartPole-v1    & $78.0 \pm 12.3$     & $485.13 \pm 8.2$   & +521\% & 97\% \\
LunarLander-v2 & $52.3 \pm 18.7$     & $241.7 \pm 15.3$   & +362\% & 96\% \\
Pendulum-v1    & $-847.2 \pm 105.4$  & $-185.4 \pm 22.1$  & +78\%  & 88\% \\
\bottomrule
\end{tabular}
\end{table}

\noindent SecFedRL consistently improves performance across discrete, sparse-reward,
and continuous control settings, validating the robustness and applicability of the
proposed policy-space detection framework.

\begin{figure}[htbp]
  \centering
  \includegraphics[width=0.95\linewidth]{image_05.png}
  \caption{Task success rate across environments under Byzantine attacks.}
  \label{fig:success}
\end{figure}

Figure~\ref{fig:success} shows the per-round CartPole-v1 task success rate---defined
as the fraction of evaluation episodes achieving $\geq$195 balanced steps (CartPole's
official solved criterion)---under 20\% Byzantine attack across 50 training rounds.
The rate is variable during the warm-up phase (rounds 0--25), reflecting the period
before behavioral baselines are fully established. The 97\% CartPole success rate
reported in Table~\ref{tab:cartpole_perf} is the aggregate performance measured across
the full 100-round converged phase of the main experiments (not visible in this 50-round
excerpt). Cross-environment success rates for LunarLander-v2 (96\%) and
Pendulum-v1 (88\%) are reported in Table~\ref{tab:multienv} and follow a similar
trajectory pattern.

\subsection{Catastrophic Collapse Prevention}

SecFedRL has implemented a multi-tiered checkpoint recovery system to protect against
the risk of irreversible training collapse. The ability of each method to protect
against training collapse is shown in Table~\ref{tab:collapse}.

\begin{table}[htbp]
\centering
\caption{Collapse Prevention Performance.}
\label{tab:collapse}
\begin{tabular}{lccc}
\toprule
\textbf{Environment} & \textbf{Collapse Events} & \textbf{Recovery Success} &
\textbf{Recovery Time (Rounds)} \\
\midrule
CartPole    & 18/100 & 100\% & 0.8 \\
LunarLander & 24/100 & 100\% & 0.9 \\
Pendulum    & 15/100 & 100\% & 0.7 \\
\bottomrule
\end{tabular}
\end{table}

\smallskip\noindent\textit{\textbf{Scope note:} The 100\% recovery success rate
reflects conditions tested: 20\% Byzantine fraction under UMP/TMP attacks across
5~seeds~$\times$~100~rounds per environment. Recovery interpolates from a saved
checkpoint (Eq.~\ref{eq:recovery}) when a reward drop exceeding 20\% is detected. The
mechanism assumes checkpoint integrity---it is not designed to defend against
\emph{targeted checkpoint-poisoning attacks}, in which an adversary corrupts the
checkpoint itself prior to a triggered recovery. Under the UMP/TMP attack models in
this study, no checkpoint-poisoning attempts were made and checkpoint integrity was
preserved throughout all runs. Additionally, recovery is not guaranteed if the collapse
threshold ($\Delta r_t > 20\%$) is not triggered---a silent performance degradation
that never crosses the threshold would not activate recovery. Both scenarios represent
acknowledged limitations of the collapse prevention mechanism and are directions for
future work.}

\begin{figure}[htbp]
  \centering
  \includegraphics[width=0.95\linewidth]{image_03.png}
  \caption{Gradient norm and reward resilience analysis under adversarial training
    conditions: (a) overall gradient norm comparison with and without Byzantine attacks,
    showing an average attack impact of +0.092 across training rounds; (b) per-layer
    gradient analysis for actor (+0.071) and critic (+0.037) components; (c) cumulative
    reward resilience demonstrating SecFedRL's recovery to near-optimal performance under
    attack; and (d) converged-phase stability comparison via gradient norm distributions
    over the last 20 rounds (std: 0.0344 without attack, 0.0462 with attack).}
  \label{fig:gradient}
\end{figure}

As shown in Figure~\ref{fig:gradient}, gradient norms remain bounded throughout
training despite Byzantine attacks. Figure~\ref{fig:gradient}(a) confirms an average
gradient norm attack impact of +0.092 across training rounds; per-layer analysis
(Figure~\ref{fig:gradient}(b)) shows the actor layer absorbs a larger share (+0.071)
than the critic (+0.037), consistent with Byzantine agents targeting the policy head.
Panel (c) shows the cumulative reward resilience: even under attack, SecFedRL recovers
toward near-optimal rewards (close to 500) while the undefended baseline cannot.
Panel (d) confirms that in the converged phase (last 20 rounds), gradient norm
standard deviation under attack (0.0462) remains only marginally higher than without
attack (0.0344), confirming stable learning dynamics. In the case of CartPole, a
representative collapse produced a major loss in episodic reward from 460.3 to just
9.4, resulting in Tier-1 recovery and restoring the episodic reward back to 483.4
after just one round. In contrast, during a collapse of a baseline method, the method
produced a permanent failure.

\subsection{Robustness to Increasing Byzantine Fractions}
\label{sec:byzantine_frac}

Robustness was further assessed by varying the proportion of Byzantine agents.
Table~\ref{tab:byzantine_frac} evaluates SecFedRL's robustness as Byzantine
participation increases in the CartPole-v1 environment. The framework maintains
near-optimal rewards and high task success rates up to 20\% adversarial agents with
98\% average per-round detection rate. Performance declines gradually as Byzantine presence grows,
demonstrating stable learning behavior. Even with 40\% Byzantine participation,
SecFedRL sustains high rewards and 90\% task success, showing strong resilience.
Unlike parameter-space baselines, which falter beyond 30\%, SecFedRL exhibits graceful
degradation under adversarial conditions.

\begin{table}[htbp]
\centering
\caption{Performance across Different Byzantine Fractions (CartPole-v1). Detection Rate is the average per-round TPR during the converged phase (rounds 30--100); compare with cumulative detection in Table~\ref{tab:detection}.}
\label{tab:byzantine_frac}
\begin{tabular}{lccc}
\toprule
\textbf{Byzantine \%} & \textbf{Final Reward} & \textbf{Avg.\ Per-Round Detection Rate (TPR)} &
\textbf{Task Success} \\
\midrule
0\%  & $493.5 \pm 5.2$    & N/A & 99\% \\
10\% & $489.2 \pm 7.1$    & 97\% & 98\% \\
20\% & $485.13 \pm 8.2$   & 98\% & 97\% \\
30\% & $472.8 \pm 11.7$   & 98\% & 95\% \\
40\% & $451.3 \pm 15.2$   & 95\% & 90\% \\
\bottomrule
\end{tabular}
\end{table}

\noindent\textbf{Threshold-corruption analysis.}
The TPR degradation at 40\% Byzantine fraction (from 98\% to 95\%) warrants further
explanation. At this fraction, the consensus threshold ($\geq 2n/3$) requires that at
least two-thirds of agents agree on a round's validity. With 40\% Byzantine agents, the
honest majority ($60\%$) still exceeds the two-thirds threshold ($\approx 66.7\%$) only
marginally. This tight margin means that Byzantine agents can cast coordinated dissent
votes to slow consensus, increasing the probability that individual rounds slip through
undetected by the behavioral analysis before the full 10-round sliding window
accumulates sufficient statistics. The result is a modest but measurable reduction in
per-round TPR. This phenomenon confirms that the safety bound of SecFedRL's consensus
mechanism aligns with its stated assumption of $f < 0.4n$ Byzantine participants: beyond
this fraction the statistical margins narrow, and the detection guarantees weaken.

\subsection{Scalability Analysis}

Table~\ref{tab:scalability} evaluates the scalability of SecFedRL as the number of
participating clients increases under 20\% Byzantine attacks. The framework maintains
high final rewards and strong detection performance even as clients scale from 10 to
50. Although computational overhead increases moderately, learning performance degrades
only slightly, indicating efficient scalability. These results demonstrate that
SecFedRL sustains robustness and stable aggregation in large-scale federated
environments.

\begin{table}[htbp]
\centering
\caption{Scalability Analysis under 20\% Byzantine Participation.}
\label{tab:scalability}
\begin{tabular}{lccc}
\toprule
\textbf{Clients} & \textbf{Final Reward} & \textbf{Detection Rate} & \textbf{Overhead} \\
\midrule
10 & $485.13 \pm 8.2$   & 98\% & 4.2\% \\
20 & $478.9 \pm 9.5$    & 96\% & 5.8\% \\
50 & $462.1 \pm 12.3$   & 95\% & 8.3\% \\
\bottomrule
\end{tabular}
\end{table}

\begin{figure}[htbp]
  \centering
  \includegraphics[width=0.95\linewidth]{image_08.png}
  \caption{Convergence of actor, critic, and total loss components.}
  \label{fig:loss}
\end{figure}

Figure~\ref{fig:loss} shows four loss component trajectories over 50 training rounds.
The actor loss (top-left) exhibits early variability before converging to a stable
negative level by round 25, consistent with the warm-up period required to establish
behavioral baselines. The critic loss (top-right) decreases steadily from $\approx$1.4
to $\approx$0.35, tracking the target convergence curve closely. The entropy loss
(bottom-left) passes through three labelled phases: Exploration (rounds 0--8),
Balanced (rounds 8--28), and Exploitation (rounds 28--50), reflecting the expected
transition from broad policy search toward focused execution. The total loss
(bottom-right) declines from $\approx$0.6 to below 0.2 by round 50, with 5-round
smoothing confirming sustained downward convergence. Final per-component stability
over the last 10 rounds: Actor = 0.019, Critic = 0.037, Entropy = 0.015, confirming
all loss components are stable post-convergence.

\subsection{Ablation Study}

A series of removal tests were carried out to evaluate how much each defensive aspect
contributed individually.

\begin{table}[htbp]
\centering
\caption{Component Ablation Evaluation under 20\% Byzantine Conditions.}
\label{tab:ablation}
\begin{tabular}{lccc}
\toprule
\textbf{Configuration} & \textbf{Final Reward} & \textbf{Detection Rate} &
\textbf{Collapses Prevented} \\
\midrule
No Defense               & $78.0 \pm 12.3$    & 0\%  & 0\%   \\
Plugin Only              & $342.7 \pm 32.5$   & 85\% & 65\%  \\
Reputation Only          & $298.5 \pm 41.2$   & 70\% & 40\%  \\
Collapse Prevention Only & $156.3 \pm 58.9$   & 0\%  & 95\%  \\
Plugin + Reputation      & $421.8 \pm 15.7$   & 95\% & 85\%  \\
Full SecFedRL            & $485.13 \pm 8.2$   & 98\% & 100\% \\
\bottomrule
\end{tabular}
\end{table}

The behavioral detection component establishes the initial adversarial signal;
reputation scoring provides persistent peer trustworthiness tracking; and collapse
prevention acts as a safety net when detection alone is insufficient. When all three
components are fully integrated, the system achieves the optimal values reported in
Table~\ref{tab:ablation}.

\textbf{Hyperparameter Sensitivity.} The number of P2P peers $k=4$ was selected
empirically: values of $k=2$ reduced diversity of behavioral cross-checks (detection
rate dropped to $\approx$91\%), while $k=8$ added overhead without measurable accuracy
improvement (detection rate stable at 98\%). The relative consensus threshold
$\delta = 5\%$ is robust to small perturbations---$\delta$ values in the range 3--8\%
produced equivalent consensus decisions across all three environments, while
$\delta < 2\%$ caused spurious rejections during the warm-up phase and $\delta > 10\%$
delayed consensus acceptance under moderate attack intensity. The entropy coefficient
$c_e = 0.01$ follows the standard A2C default and was not swept in this study;
sensitivity to $c_e$ is left as future work.

\subsection{Statistical Significance}

Statistical significance was assessed using independent two-sample $t$-tests at
$\alpha = 0.001$ on per-round reward observations drawn from the converged phase of
training. Each method was evaluated over 5 independent random seeds $\times$ 3
environments $\times$ 100 training rounds = 1,500 total round-reward observations per
method. Pairwise $t$-tests compare reward distributions sampled from the final 30
rounds (rounds 71--100) of each trial---5 seeds $\times$ 3 environments $\times$ 30
rounds = 450 observations per method---yielding $t(898)$ degrees of freedom for each
comparison.

\textbf{Note on cross-environment pooling:} CartPole-v1 (rewards $\approx$ +400--500),
LunarLander-v2 (rewards $\approx$ +52--241), and Pendulum-v1 (rewards $\approx$ $-185$
to $-847$) have different reward scales and sign conventions. Pooling observations
across environments amplifies the apparent effect size for the FedAvg comparison
(which performs very poorly across all three environments). Per-environment effect
directions are consistent: SecFedRL strictly dominates all baselines in all three
environments across all seeds. Pooled statistics are reported here as a compact
summary; environment-stratified analyses are provided within
Tables~\ref{tab:cartpole_perf}--\ref{tab:byzantine_frac}.

All pairwise comparisons reach significance at $p < 0.001$:
\begin{itemize}
  \item SecFedRL vs.\ FedAvg:         $t(898) = 47.23$, $p < 0.001$, $d = 3.15$
  \item SecFedRL vs.\ Krum:           $t(898) = 12.45$, $p < 0.001$, $d = 0.83$
  \item SecFedRL vs.\ Trimmed-Mean:   $t(898) = 14.67$, $p < 0.001$, $d = 0.98$
\end{itemize}

The large effect size for FedAvg ($d = 3.15$) reflects the fundamental performance
difference between a defended and an undefended system under adversarial conditions.
The moderate-to-large effect sizes vs.\ Krum and Trimmed Mean ($d = 0.83$, $0.98$)
confirm practical superiority beyond statistical significance.

\noindent\textbf{Per-environment consistency.}
Effect directions are consistent across all three environments: SecFedRL strictly
dominates all baselines in CartPole-v1, LunarLander-v2, and Pendulum-v1 across all
five seeds. Environment-stratified effect sizes (available in the full experimental
data at \url{https://github.com/Yog37h/byzsecure}) range from $d = 0.71$
(SecFedRL vs.\ Krum in LunarLander-v2) to $d = 4.12$ (SecFedRL vs.\ FedAvg in
CartPole-v1), confirming that the aggregate effect is not driven by any single
environment.

\subsection{Discussion}

The Figures~\ref{fig:convergence} through~\ref{fig:loss}, in combination, summarize
the findings and describe: manner and speed of convergence, detection rate for
Byzantine actors, success rates in completing any task in which all parties shared a
common goal, stability of actors in either direction (toward or away from policy
established), and scalability.

The proposed SecFedRL framework demonstrates strong resistance to Byzantine adversaries
by combining policy-space behavioral analysis, reputation-based peer weighting, and
checkpoint-based recovery---without requiring a central coordinator or ground-truth
labels. SecFedRL achieves a 98\% average per-round detection rate with near-zero false
positives, and shows faster convergence under adversarial conditions than
parameter-space baselines in CartPole-v1 (5 rounds vs.\ $>$50 rounds for Krum/Trimmed
Mean to reach reward threshold $\geq$450 in the poisoned setting; see
Figure~\ref{fig:convergence} and Table~\ref{tab:cartpole_perf}).

Over the course of 15 independent experiment runs (5 seeds $\times$ 3 environments),
SecFedRL improved episodic rewards by 78--521\% over the undefended FedAvg baseline,
with environment-specific improvements of +521\% on CartPole-v1, +362\% on
LunarLander-v2, and +78\% absolute reward improvement on Pendulum-v1 (a
continuous-control environment with negative-valued rewards where multiplicative
comparisons are not directly applicable), while consistently resulting in 88--97\%
task success rates across all experimental conditions (CartPole: 97\%, LunarLander:
96\%, Pendulum: 88\%). Although the SecFedRL framework requires a brief warm-up period
to establish the behavioral baselines that it will operate on, it is also very
well-suited for real-world applications of decentralized learning such as multi-robot
coordination, autonomous fleet management, and federated IoT systems. All of the work
described in this paper establishes the necessary scalable and robust basis to enable
secure P2P FRL.

\noindent\textbf{Limitations.}
Several limitations bound the current results. \textit{Convergence guarantees:}
the experiments demonstrate empirical convergence across five seeds and three
environments but do not establish formal theoretical convergence guarantees for the
joint reputation-aggregation-consensus loop under Byzantine conditions. \textit{Environment
homogeneity:} the virtual trajectory mechanism requires all agents to share the same
environment instance (see Section~\ref{sec:threat}); heterogeneous multi-task settings
are not addressed. \textit{Advanced attacks:} the detection framework is evaluated
against UMP and TMP attacks; mimicry attacks---where adversaries calibrate updates to
remain within the benign behavioral distribution---and checkpoint-poisoning attacks
remain open vulnerabilities. These are priority directions for future work.

%% ════════════════════════════════════════════════════════════════════════
%% §6  CONCLUSION
%% ════════════════════════════════════════════════════════════════════════
\section{Conclusion}
\label{sec:conclusion}

In this paper, SecFedRL was presented as a comprehensive framework for P2P FRL that is
Byzantine resilient. The approach that was taken to develop SecFedRL includes four key
components: (1) Policy Space Behavioral Detection; (2) Ground Truth Free Reputation
Learning; (3) Catastrophic Collapse Prevention; and (4) Lightweight and Decentralized
Consensus. These four components combined to provide a solution to the
security limitations present within the current parameter space defence mechanisms of
reinforcement learning.

SecFedRL provides excellent performance with regard to Byzantine detection (average
per-round True Positive Rate of 98\% at 20--30\% Byzantine fraction, with 100\%
cumulative detection within 25--45 rounds and graceful degradation to 95\% average
per-round TPR at 40\%, and near-zero False Positive Rate) across the CartPole-v1,
LunarLander-v2, and Pendulum-v1 environments using a seven metric-based policy space
behavioral evaluation framework. The reputation mechanism utilized within SecFedRL
distinguishes malicious vs.\ benign clients (Byzantine clients obtain reputation
scores 0.42--0.73; benign participants 0.64--0.89). While raw reputation scores
partially overlap in the 0.64--0.73 boundary region, the composite suspicion score
$S_j$ provides a separating margin of $0.15+$ standard deviations above the benign
mean, enabling correct classification without any dependencies on trusted validation
data sets. In addition, the three-tiered checkpoint recovery mechanisms of SecFedRL
prevent irreversible policy collapse through timely checkpoint recovery.

Through comprehensive testing across 15 independent experiment runs (5 seeds $\times$
3 environments, 100 rounds each), SecFedRL shows 362--521\% improvement in final
episodic rewards over the undefended baseline across discrete-action environments
(CartPole-v1: +521\%, LunarLander-v2: +362\%), with significant improvement also in
the negative-reward Pendulum-v1 setting (+78\% absolute reward improvement), and
88--97\% task success rates across environments (CartPole: 97\%, LunarLander: 96\%,
Pendulum: 88\%) while under significant Byzantine attacks. There is graceful
degradation of performance with respect to up to 40\% of agents acting adversarially
during the task.

Additionally, SecFedRL demonstrates faster convergence under adversarial conditions than leading
state-of-the-art methods (5 rounds to reach high-performance reward threshold $\geq$450
in CartPole-v1 under poisoning, compared to more than 50 rounds for parameter-space
baselines; see Figure~\ref{fig:convergence} and Table~\ref{tab:cartpole_perf}), while
only requiring 4.2\% additional computational time, indicating the practicality of SecFedRL for
real-world deployment of decentralized learning environments. In general, SecFedRL
provides a method for enabling trustworthy learning through collaborative communication
between the various agents in an adversarial, decentralized environment. This opens up
many different types of applications to use SecFedRL such as coordination of multiple
robots, fleets of autonomous vehicles, and federated Internet of Things (IoT)
ecosystems. This work establishes a strong foundation for secure, scalable, and
resilient federated reinforcement learning, advancing the development of decentralized
autonomous intelligent systems.

Future work will explore extending SecFedRL to high-dimensional visual
reinforcement learning applications, establishing formal convergence guarantees
for the joint reputation-aggregation-consensus loop among all agents, providing
adaptive defences against colluding adversaries via a decentralized mechanism,
and establishing an asynchronous communication protocol for heterogeneous edge
platform agents.

%% ════════════════════════════════════════════════════════════════════════
%% DATA AVAILABILITY STATEMENT
%% ════════════════════════════════════════════════════════════════════════
\section*{Data Availability Statement}

The implementation of SecFedRL, including training scripts, evaluation code, and
experimental configurations for CartPole-v1, LunarLander-v2, and Pendulum-v1
environments, is publicly available at \url{https://github.com/Yog37h/byzsecure}.
All experiments use OpenAI Gym v0.26.x with PyTorch v2.0. The five independent runs
reported in Section~\ref{sec:results} were conducted using random seeds $\{0, 1, 2, 3, 4\}$.
Hyperparameter configurations are listed in Table~\ref{tab:hyperparams}. The number
of local training steps per communication round per agent ($T_{\text{local}}$) is
128--512 steps depending on environment complexity; full configurations are available
in the repository.

%% ════════════════════════════════════════════════════════════════════════
%% REFERENCES
%% ════════════════════════════════════════════════════════════════════════
\begin{thebibliography}{99}

\bibitem{ref1}
M.~Karami, F.~Ghassemi, H.~Kebriaei, and H.~Azadegan,
``OptiGradTrust: Byzantine-robust federated learning with multi-feature gradient
analysis and reinforcement learning-based trust weighting,''
\textit{arXiv preprint arXiv:2507.23638}, 2025.
\url{https://arxiv.org/abs/2507.23638}

\bibitem{ref2}
P.~Kairouz and H.~B.~McMahan,
``Advances and open problems in federated learning,''
\textit{Foundations and Trends in Machine Learning}, vol.~14, no.~1--2, pp.~1--210,
2021.
\url{https://doi.org/10.1561/2200000083}

\bibitem{ref3}
H.~Zhu, J.~Xu, S.~Liu, and Y.~Jin,
``Federated learning on non-IID data: A survey,''
\textit{Neurocomputing}, vol.~465, pp.~371--390, 2021.
\url{https://doi.org/10.1016/j.neucom.2021.07.098}

\bibitem{ref4}
K.~Bonawitz, V.~Ivanov, B.~Kreuter, A.~Marcedone, H.~B.~McMahan, S.~Patel,
D.~Ramage, A.~Segal, and K.~Seth,
``Practical secure aggregation for privacy-preserving machine learning,''
in \textit{Proc.\ 2017 ACM SIGSAC Conference on Computer and Communications
Security (CCS)}, pp.~1175--1191, 2017.
\url{https://doi.org/10.1145/3133956.3133982}

\bibitem{ref5}
V.~Shejwalkar, A.~Houmansadr, P.~Kairouz, and D.~Ramage,
``Back to the drawing board: A critical evaluation of poisoning attacks on production
federated learning,''
in \textit{Proc.\ 2022 IEEE Symposium on Security and Privacy (SP)},
pp.~1354--1371. IEEE, 2022.
\url{https://doi.org/10.1109/SP46214.2022.9833647}

\bibitem{ref6}
A.~Sokhankhosh and S.~Rouhani,
``Proof-of-Collaborative-Learning: A multi-winner federated learning consensus
algorithm,''
in \textit{Proc.\ 2024 IEEE International Conference on Blockchain},
pp.~370--377. IEEE, 2024.
\url{https://doi.org/10.1109/blockchain62396.2024.00055}

\bibitem{ref7}
Y.~Li and X.~Lyu,
``Convergence analysis of sequential federated learning on heterogeneous data,''
\textit{Advances in Neural Information Processing Systems}, vol.~36,
pp.~56700--56755, 2023.
\url{https://proceedings.neurips.cc}

\bibitem{ref8}
V.~Shejwalkar and A.~Houmansadr,
``Manipulating the Byzantine: Optimizing model poisoning attacks and defenses for
federated learning,''
in \textit{Proc.\ Network and Distributed System Security Symposium (NDSS)}, 2021.
\url{https://doi.org/10.14722/ndss.2021.24498}

\bibitem{ref9}
G.~Lan, X.-Y.~Liu, Y.~Zhang, and X.~Wang,
``Communication-efficient federated learning for resource-constrained edge devices,''
\textit{IEEE Transactions on Machine Learning in Communications and Networking},
vol.~1, pp.~210--224, 2023.
\url{https://doi.org/10.1109/TMLCN.2023.3309773}

\bibitem{ref10}
J.~Xu, Z.~Zhang, and R.~Hu,
``Achieving Byzantine-resilient federated learning via layer-adaptive sparsified model
aggregation,''
in \textit{Proc.\ 2025 IEEE/CVF Winter Conference on Applications of Computer Vision
(WACV)}, pp.~1508--1517. IEEE, 2025.
\url{https://doi.org/10.1109/wacv61041.2025.00154}

\bibitem{ref11}
T.~Lin, L.~Kong, S.~U.~Stich, and M.~Jaggi,
``Ensemble distillation for robust model fusion in federated learning,''
\textit{arXiv preprint arXiv:2006.07242}, 2020.
\url{https://arxiv.org/abs/2006.07242}

\bibitem{ref12}
Y.~Du, S.~Yang, and K.~Huang,
``High-dimensional stochastic gradient quantization for communication-efficient edge
learning,''
in \textit{Proc.\ 2019 IEEE Global Conference on Signal and Information Processing
(GlobalSIP)}, pp.~1--5. IEEE, 2019.
\url{https://doi.org/10.1109/GlobalSIP45357.2019.8969082}

\bibitem{ref13}
B.~Tegin and T.~M.~Duman,
``Federated learning with over-the-air aggregation over time-varying channels,''
\textit{IEEE Transactions on Wireless Communications}, vol.~22, no.~8,
pp.~5671--5684, 2023.
\url{https://doi.org/10.1109/TWC.2023.3235894}

\bibitem{ref14}
B.~Deressa and M.~A.~Hasan,
``TrustBandit: Optimizing client selection for robust federated learning against
poisoning attacks,''
in \textit{Proc.\ IEEE INFOCOM 2024 Workshops}, pp.~1--8. IEEE, 2024.
\url{https://doi.org/10.1109/INFOCOMWKSHPS61880.2024.10620802}

\bibitem{ref15}
F.~X.~Fan, Y.~Ma, Z.~Dai, W.~Jing, C.~Tan, and B.~K.~H.~Low,
``Fault-tolerant federated reinforcement learning with theoretical guarantee,''
\textit{Advances in Neural Information Processing Systems}, vol.~34, 2021.
\url{https://arxiv.org/abs/2110.14074}

\bibitem{ref16}
M.~Fang, X.~Cao, J.~Jia, and N.~Z.~Gong,
``Local model poisoning attacks to Byzantine-robust federated learning,''
in \textit{Proc.\ 29th USENIX Security Symposium}, pp.~1605--1622, 2020.
\url{https://arxiv.org/abs/1911.11815}

\bibitem{ref17}
Z.~Xie and S.~Song,
``FedKL: Tackling data heterogeneity in federated reinforcement learning by penalizing
KL divergence,''
\textit{IEEE Journal on Selected Areas in Communications}, vol.~41, no.~4,
pp.~1227--1242, 2023.
\url{https://doi.org/10.1109/JSAC.2023.3242734}

\bibitem{ref18}
L.~Ni, X.~Gong, J.~Li, Y.~Tang, Z.~Luan, and J.~Zhang,
``rFedFW: Secure and trustable aggregation scheme for Byzantine-robust federated
learning in Internet of Things,''
\textit{Information Sciences}, vol.~653, p.~119784, 2024.
\url{https://doi.org/10.1016/j.ins.2023.119784}

\bibitem{ref19}
P.~Blanchard, E.~M.~El~Mhamdi, R.~Guerraoui, and J.~Stainer,
``Machine learning with adversaries: Byzantine-tolerant gradient descent,''
\textit{Advances in Neural Information Processing Systems}, vol.~30, 2017.
\url{https://papers.nips.cc/paper_files/paper/2017/hash/f4b9ec30ad9f68f89b29639786cb62ef-Abstract.html}

\bibitem{ref20}
Y.~Chen, L.~Su, and J.~Xu,
``Distributed statistical machine learning in adversarial settings: Byzantine gradient
descent,''
\textit{Proc.\ ACM on Measurement and Analysis of Computing Systems}, vol.~1, no.~2,
pp.~1--25, 2017.
\url{https://doi.org/10.1145/3154503}

\bibitem{ref21}
C.~Xie, O.~Koyejo, and I.~Gupta,
``Generalized Byzantine-tolerant SGD,''
\textit{arXiv preprint arXiv:1802.10116}, 2018.
\url{https://arxiv.org/abs/1802.10116}

\bibitem{ref22}
D.~Yin, Y.~Chen, K.~Ramchandran, and P.~Bartlett,
``Byzantine-robust distributed learning: Towards optimal statistical rates,''
in \textit{Proc.\ 35th International Conference on Machine Learning (ICML)},
PMLR vol.~80, 2018.
\url{https://arxiv.org/abs/1803.01498}

\bibitem{ref23}
L.~Lyu, H.~Yu, J.~Zhao, and Q.~Yang,
``Threats to federated learning,''
in \textit{Federated Learning}, Lecture Notes in Computer Science, vol.~12500,
pp.~3--16. Springer, 2020.
\url{https://doi.org/10.1007/978-3-030-63076-8_1}

\bibitem{ref24}
Q.~Yang, Y.~Liu, T.~Chen, and Y.~Tong,
``Federated machine learning: Concept and applications,''
\textit{ACM Transactions on Intelligent Systems and Technology}, vol.~10, no.~2,
pp.~1--19, 2019.
\url{https://doi.org/10.1145/3298981}

\bibitem{ref25}
H.~B.~McMahan, E.~Moore, D.~Ramage, S.~Hampson, and B.~Ag\"{u}era~y~Arcas,
``Communication-efficient learning of deep networks from decentralized data,''
in \textit{Proc.\ 20th International Conference on Artificial Intelligence and
Statistics (AISTATS)}, PMLR vol.~54, 2017.
\url{https://arxiv.org/abs/1602.05629}

\bibitem{ref26}
K.~Cui, X.~Feng, L.~Wang, and Z.~Liu,
``Dual-verifiable federated learning with vector commitments against collusion
attacks,''
\textit{IEEE Transactions on Dependable and Secure Computing}, pp.~1--17, 2026.
\url{https://doi.org/10.1109/TDSC.2026.3674195}

\bibitem{ref27}
H.~Chen, X.~Chen, L.~Peng, and R.~Ma,
``FLRAM: Robust aggregation technique for defense against Byzantine poisoning attacks
in federated learning,''
\textit{Electronics}, vol.~12, no.~21, p.~4463, 2023.
\url{https://doi.org/10.3390/electronics12214463}

\bibitem{ref28}
J.~Hou, X.~Liang, N.~Gui, L.~Wang, and S.~Zheng,
``Byzantine-robust decentralized federated learning based on feature extraction layer
measuring and delayed aggregation,''
\textit{Knowledge-Based Systems}, vol.~337, p.~115392, 2026.
\url{https://doi.org/10.1016/j.knosys.2026.115392}

\bibitem{ref29}
Y.~Xia, W.~Yu, and Q.~Li,
``Byzantine-resilient federated learning via distributed optimization,''
in \textit{Proc.\ 2025 33rd European Signal Processing Conference (EUSIPCO)},
pp.~2672--2676. IEEE, 2025.
\url{https://doi.org/10.23919/eusipco63237.2025.11226343}

\bibitem{ref30}
S.~Salgia and Y.~Chi,
``The sample-communication complexity trade-off in federated Q-learning,''
\textit{Advances in Neural Information Processing Systems}, vol.~37,
pp.~39694--39747, 2024.
\url{https://proceedings.neurips.cc}

\end{thebibliography}

\end{document}
